\chapter[Sermon 10. The First Part of the Third Epistle of the Constancy and Confession of Christ in the Time of Persecution.]{The First Part of the Third Epistle of the Constancy and Confession of Christ in the Time of Persecution.}
\label{sermon:010}
\markright{Sermon 10. The First Part of the Third Epistle of the Constancy ...}

And to the messenger of the congregation in Pergamum write: This says he who has the sharp sword with two edges. I know your works and where you dwell, even where Satan’s throne is, and you hold fast to my name and have not denied my faith. And that in the days in which Antipas was my faithful witness, who was slain among you where Satan dwells.

\index{Repentance}The third epistle among those seven heavenly ones proceeding from the right hand of God, is written to the pastor and congregation of Pergamum. The argument is as follows: First, he commends the constancy of their faith in cruel persecutions. Then, he rebukes those who clung to the sect of the Nicolaitans. Afterward, he exhorts them to repentance. And this doctrine he applies afterward to all congregations throughout the world. Finally, he promises most ample rewards to the faithful. We understand from this that the congregation of Pergamum is presented as a type or a mirror to all churches, showing how they ought to walk before the Lord: first, whenever persecution arises; secondly, when heresies break out. For by its example, he teaches all to suffer adversity patiently and openly to profess the true faith, and also by Scripture to reprove heresies, and in fleeing from them to despise them.

However, all the Epistles have certain things in common, and especially three. It plainly expresses to whom the Epistle is sent, as in this present case, to the messenger of the congregation of Pergamum, namely to the pastor, whoever he was (perhaps Antipas), and to the whole congregation, as was said before. It is shown moreover who speaks here, or who is the author of this Epistle: even the Lord himself. This lends authority to the writing. For it is not to be thought that the word of God is not as it is spoken because it is written by a man, dictated by a man, or written with ink, either on paper or parchment. For these things do not make the word of God cease to be the word of God, any more than water ceases to be water if it runs out of a conduit of wood, lead, brass, or stone. For water always remains water. The diversion of the conduit pipes does not change it from being water, as it is in substance. So says Paul, that he may be bound, but the word of God is not bound. A man may be stoned, hanged, or burned, being a preacher of God’s word; the word of God that was put in the mouth of the preacher is not burned. The Lord puts it in the mouth of another, that the truth should not be extinguished, but continually might sound in the church. Finally, without cause, at the beginning of every Epistle, Christ intimates that he knows all things of the church. I said before that this is as it were the foundation of the fear of God, and of his true service. For imagine a man who thinks that God neither sees what men do, nor knows what they think in their hearts. Shall not that man fall into all ungodliness? He will cry, let us do what we list, since God knows not what we do. Who will not cast off the hope of reward, and the practice of good works, after he is once persuaded that God knows not our works? But if he knew them not, how would he judge the world?

Nevertheless, in every epistle, be certain of special and peculiar things. Of this sort, in the epistle of Pergamum, is that from the first vision and description of Christ, at the beginning of the epistle, he takes up the sworn, sharp, and two-edged sword, which we heard come out of the mouth of Christ. By this is signified the judicial power of equity and justice, and also the deliverance of the good and the punishment of the evil, for the sword is given to the magistrate as an authority to punish the evil and defend the good. Christ himself defends his own, and he hews his adversaries in pieces. The sword is the very word of God, most sharp, two-edged, and piercing the hearts, for it animates the godly and discourages the wicked. Christ therefore governs his Church as a Judge and defender, most rightful and just, which has his sword not in his hands, but in his mouth, and with his spirit and word comforts and preserves the faithful, but fears and brings death to the unbelievers. Full rightly therefore is this beginning applied to the cause that follows, concerning the cross of the faithful and expelling the Nicolaitans. For it is Christ, by whose virtue these things are successfully brought to pass.

Furthermore, the specific actions of this congregation are recorded. He commends the church for its remarkable steadfastness in faith and its public declaration of that faith, even amidst perilous trials, temptations, and persecutions. It appears to be a straightforward account showing that the Lord is aware of their suffering and how severely they are afflicted, yet praise is interwoven throughout. This commendation serves as an encouragement, urging them to continue in the actions they have already undertaken.

He says he is not ignorant of where the church of Pergamum dwells: even there, truly, where Satan has fixed his seat or throne. That is to say, I know in what case you are, in what dangers, and with whom you are contending. He says not, I know that you sit in the seat of Satan, but, I know that you dwell there where Satan has his seat. Christ therefore is not ignorant of the labors, sorrows, and temptations of the faithful. And the knowledge of Christ has a certain peculiar quality. For Christ so knows the matters of the faithful that he is both touched by them and has also a consideration or respect for his servants. And we see how Christ also places his throne there where the Devil has his seat, just beside it. At length he thrusts him out of his seat.

\index{Arethas of Caesarea}\index{Rome}And for two causes Pergamos seems to be called the seat, throne, and kingdom of the devil. For first as Arethas has admonished, in superstition and worshipping of idols it excelled all Asia, which nevertheless was most corrupt. Pergamos was the most ancient and famous city of Asia or of Mysia and Phrygia, renowned by king Attalus and Eumenes. For the same was the princely palace of king Attalus, which came into the hands of the Romans by the legacy of kings, who were most addicted to idolatry. Strabo speaks much hereof in the 13th book. Moreover this place was also, as Pliny seems to signify in the 5th book, the 30th Chapter, most noble and frequented, by reason the lieutenant or governor there inhabited, who at the commandment of the emperor Domitian, persecuted the true faith of Christ, imprisoning, scourging and afflicting all that professed Christ. By good reason therefore is Pergamos called the seat or throne of the devil. For he is a liar, and the father of lying, and a murderer from the beginning: which the Lord also testifies in the 8th of John. For because therefore at Pergamos reigned heathens, lying, idolatry, superstition, the oppression and murder of good men, it is rightly called the seat or throne of the devil. This appears to be a slander not to be dissembled or suffered. For Rome seemed to herself established forever, and which the gods favored, who had sent them victory over most great nations, and given the empire of the whole world: in the which city justice and religion might seem to be observed. And therefore that this seat of justice and religion should be called the seat of Satan, might be thought both blasphemy and treason. But this doth the only begotten son of God from the right hand of his father pronounce against Rome, against Pergamos, and against all the supporters of Rome. Who shall accuse him of temerity, of rashness, or of bitter speaking? Light persons are doubtless angry, and very sharp wits will be offended, in case they be called by their own names, and be called as they are in deed. For such is the glory of virtue, that all men covet the same even the open enemies of virtue, so that no man will seem to be void of virtue: and such is the corruption and darkness of man’s mind, that he would be that he is not, and would not be that he is. Thereof comes all this impatience in the whole world: when a mattock is called a mattock, and a fig a fig, as the proverb is. Is an harlot therefore no harlot, because she will not be called an harlot? Yes verily is she an harlot, and a shameful harlot, and though she deny never so oft that she is a whore, yet is she an whore nevertheless, and remains a whore. So the seat or throne of Satan is at this day Rome itself, which will seem to be the seat of Christ and the seat Apostolical. For the work and instruction of the devil therein abounds. Finally all cities, towns, and places, wherein truth, godliness, religion and virtue are exiled, wherein the preaching of God’s truth and correction of most corrupt manners have no place, wherein filthiness and uncleanness, bawdy songs and not spiritual Psalms, wherein craft and deceit, surfeiting, murder, adultery, oppression of good people and of godly religion triumphs, be the seats of Satan, however they be called the most Christian and Catholic cities, and worshippers of the right and Christian faith. This thing Jesus Christ the very son of God says, cries, affirms, repeats, and even with a majesty pronounces. For by and by after the murder of Antipas, he adds: where Satan dwells. And these things are doubtless true, which Christ says and pronounces in the Church: and most false be the things which this most sinful world here alleges against the words of Christ.

But the Lord greatly praises this very thing, that in so precarious and unfortunate a place they have stood steadfastly until now, and could not be subdued in the very stronghold of Satan. Here we learn that it is lawful, as occasion serves, to dwell among a perverse nation; yet so that we are not made conformed to them in any way, either in manners or superstition. And since it is dangerous to dwell among the ungodly, as it were to touch pitch with our hands, you shall commit nothing offensive against the Lord if you go to a safer place, where there is less danger and more opportunity for all godliness. Indeed, when you may conveniently pass to such places, you risk perilously upon rocky ground, where you may chance at last to suffer shipwreck.

And he allows chiefly two things in this church, first that they hold the name of Christ. For the Greek word is not to be touched lightly, but to hold fast, so that it cannot with force be plucked away that you hold. And so they held Christ most deeply fixed in their minds. The name of Christ is the wholesome working of our redemption and sanctification, besides which there is no other name, as Saint Peter says, whereby we may be saved. They cleaved therefore unto Christ, as we read of the apostles in John 6. And it is necessary that every one of us hold fast the mystery of salvation rooted in our hearts.

Secondly, it is not enough to retain the mystery of salvation in our heart, unless we profess it also with full and open mouth. Whereupon he adds straightway, “I have not denied my faith.” Behold how he calls it faith now, which of late he called the name of Christ. And he calls it properly his faith, that is, not devised or invented by men, but set forth by Christ himself by the word of his truth. This true, right, and catholic faith must we confess and not deny, and profess it expressly as well in words as in works.

\index{Antichrist}Christ and his Gospel are denied in many ways. They are denied by silence, when we remain silent when we ought to speak primarily for the glory of God. Christ is denied through dissimulation, as when Peter says, “I do not know what you are saying.” For he knew very well what the servant girl said; but fear caused him to dissemble. He is denied when Christ and his truth are plainly and with explicit words denied. He is denied with a figurative confession, when in deed we confess something, but so darkly and so vaguely that it is unknown what we profess. He is denied when we pretend in our heart that we keep the true doctrine, and deny it in our works, by bowing before idols, attending profane churches, participating in the ceremonies of Antichrist, kneeling on the ground and worshipping that which our conscience gives us, and the faith taught by the apostles instructs us is not God. And truly, all this denial arises from fear and our corrupt affections. If there were as assuredly a reward proposed by men for confessing him, as you are greatly afraid to suffer pain if you do confess, there would seem no difficulty at all in sincerely professing Christ. Wherefore, when you deny or dissemble, you do so from fear. But such timorous and fearful deniers the Lord shuts out of his kingdom. Therefore, the world being despised, the name of the Lord must be confessed boldly and without fear, according to the doctrine of Christ, Matthew 10; Mark 8.

This declaration of faith from the congregation of Pergamum is amplified and highly commended because of the time. For it is a great thing to profess Christ not in peace, but in very troubled times. It is clear that the church of Pergamum confessed Christ in the midst of persecution, during which the holy martyr Antipas was executed. Therefore, the profession was noble. It is commonly said, but these men saw Antipas slain, and yet could not be deterred from the true faith. These things are indeed expressed briefly, but carry a profound meaning for all churches to follow. Some others read it differently in my days. But the Complutensian copy is better, which reads, “in the days wherein Antipas, \&c.” As though he should say, “And you have confessed my name in those days in which Antipas was my faithful witness, who for the same reason was also slain.”

Antipas is commended, and as it were canonized by the very Son of God. He is praised as a witness, that is, a martyr: and indeed a faithful witness, by testifying, teaching, confessing, and keeping his faith to the Lord, even to the end. Acts 13. Perhaps he was pastor of this Church, or some other man of singular constancy among the faithful. Certainly faith, not torment, makes martyrs. And because this martyr is praised by Christ, we understand that the agonies and conflicts of martyrs should be preached in the church of Christ, and many be excited and exhorted to follow their steps. Therefore we affirm that the holy martyrs of God are honored, but not to be worshipped or called upon. We condemn all those that speak against holy martyrs, and associate them with those that slew them. But concerning the worshipping of Saints, I have spoken elsewhere more at length; we learn hereof also, that they do not die forever, those who die in this world for the name of Christ: neither that the martyrs be polluted with worldly reproach, considering how they are commended by the mouth of God. To Christ therefore, king of martyrs, be honor, praise, and glory world without end. Amen.
